
\label{largepitchangle}

\subsection{Relativistic electron-electron collisions}

The starting point of the calculations for the avalanche source term in the
runaway process is close Coulomb collisions between a thermal or sub-thermal
electron and and ultrarelativistic electron. Since the realmtivistic
electron has a very high momentum, loosing a small amount of it may not
affect its dynamics significantly (it remains in the runaway phase space
after the momentum transfer), but for the other electron the gain in energy
may be so large that the drag force is suddenly strongly decreasing. For an
electron whose kinetic energy is closely below the runaway energy threshold,
such a process may make it runaway, leading by multiple interaction to the
avalanche effect.

The electron is characterized by the 4-momentum vector $\mathbf{P}$ defined
as 
\begin{equation}
\mathbf{P\equiv }\gamma \left( 1,\frac{\mathbf{v}}{c}\right) 
\end{equation}
where $\gamma =\sqrt{1+p^{2}}=1/\sqrt{1-\left( v/c\right) ^{2}}$ is the
normalized total energy $E_{T}=\gamma m_{e}c^{2},$ and $\mathbf{v}/c$ is the
normalized momentum $p=\gamma m_{e}c^{2}$. The scalar product for the
4-momentum vector $\mathbf{P}$ is defined by the relation 
\begin{equation}
\mathbf{A\cdot B=}A_{1}B_{1}-A_{2}B_{2}-A_{3}B_{3}-A_{4}B_{4}
\end{equation}
according the relativistic theory, and therefore 
\begin{equation}
\mathbf{P\cdot P=}\gamma ^{2}\left( 1-\frac{\mathbf{v\cdot v}}{c^{2}}\right)
=\gamma ^{2}\left( 1-\frac{v^{2}}{c^{2}}\right) =1
\end{equation}

Let consider two electrons, one moving at high velocity close to the speed
of light $c$, the other being at rest in the frame of mass, since its
velocity is considering to be negligible as compared to $c$. In an elastic
collision, both energy and momentum are conserved. If initial and final
4-momentum values are $(\mathbf{P}_{0}\mathbf{,}$ $\mathbf{P}_{1}\mathbf{)}$
and $(\mathbf{P}_{2}\mathbf{,}$ $\mathbf{P}_{3}\mathbf{)}$ for the fast and
slow electrons, one has the relation 
\begin{equation}
\mathbf{P}_{0}+\mathbf{P}_{1}=\mathbf{P}_{2}+\mathbf{P}_{3}
\end{equation}
or 
\begin{equation}
\mathbf{P}_{0}+\mathbf{P}_{1}-\mathbf{P}_{2}=\mathbf{P}_{3}
\end{equation}

Squaring both sides and using the definition of the scalar product for
4-momentum vectors, it comes 
\begin{equation}
1+\mathbf{P}_{0}\mathbf{\cdot P}_{1}-\mathbf{P}_{0}\mathbf{\cdot P}_{2}-%
\mathbf{P}_{1}\mathbf{\cdot P}_{2}=0
\end{equation}
and using the formulation 
\begin{equation}
\mathbf{P}_{i}\mathbf{\cdot P}_{j}=\gamma _{i}\gamma _{j}\left( 1-\frac{%
\mathbf{v}_{i}\mathbf{\cdot v}_{j}}{c^{2}}\right) 
\end{equation}
with the assumptions that $\mathbf{v}_{0}\sim 0$ and $\gamma _{0}\sim 1$
since the target electron is nearly at rest before the collision, one
obtains 
\begin{equation}
1+\gamma _{1}-\gamma _{2}-\gamma _{1}\gamma _{2}\left( 1-\frac{\mathbf{v}_{1}%
\mathbf{\cdot v}_{2}}{c^{2}}\right) =0
\end{equation}
or 
\begin{equation}
1+\gamma _{1}-\gamma _{2}-\gamma _{1}\gamma _{2}+p_{1}p_{2}\xi ^{*}=0
\label{eecoll}
\end{equation}
where $\xi ^{*}$ is the cosine of the pitch-angle between the incoming
primary electron and the scattered secondary electron as defined in Fig.
***. Here, 
\begin{eqnarray}
\xi ^{*} &=&-\frac{\left( 1+\gamma _{1}\right) \left( 1-\gamma _{2}\right) }{%
p_{1}p_{2}}  \nonumber \\
&=&-\frac{\left( 1+\sqrt{1+p_{1}^{2}}\right) \left( 1-\sqrt{1+p_{2}^{2}}%
\right) }{p_{1}p_{2}}  \nonumber \\
&=&\sqrt{\left( \frac{\gamma _{1}+1}{\gamma _{1}-1}\right) \left( \frac{%
\gamma _{2}-1}{\gamma _{2}+1}\right) }  \label{pitchanglerunaway}
\end{eqnarray}
and since $\left| \xi ^{*}\right| \leq 1,$ the validity of the relation
imposes $p_{2}\leq p_{1}.$ For very energetic runaway electrons, the
following condition  $\gamma _{1}\gg 1$ is fullfiled, and
\begin{equation}
\xi ^{*}\approx \sqrt{\frac{\gamma _{2}-1}{\gamma _{2}+1}}
\end{equation}

This relation relates the cosine of the scattering angle $\xi _{2}$ of the
slow electron to its momentum $p_{2}$ after the collisions with the fast
electron. When $p_{2}\sim 0,$ i.e. when the gain in energy is small, $\xi
^{*}$ is close to zero, which means that the scattering is perpendicular to
the direction of the runaway electron. In that case, the slow electron gains
only perpendicular energy with respect to the magnetic field line direction.
Conversely, when $p_{2}\gg 1,$ the slow electron becomes very energetic in
the collision process, and in this case $\xi ^{*}\sim 1,$ which means that
the electron has a high probability to fall into the runaway region in the
momentum phase space. It is important to notice that $\xi ^{*}>0,$ i.e. that
the rest electrons can only be scattered in the direction of the primary
runaway electron, and the locus of the end point of the scattered momentum
is along an ellipse as indicated in Ref. \cite{chi98}. Let define $p_{2\Vert
}=p_{2}\xi ^{*},$ the projection of the scattered electron momentum on the
momentum direction of the primary electron. By construction $\gamma
_{2}^{2}=1+$ $p_{2\Vert }^{2}+$ $p_{2\bot }^{2}.$ From relation \ref{eecoll}%
, one obtains
\begin{equation}
\left( 1+\gamma _{1}\right) \left( 1-\gamma _{2}\right) +p_{1}p_{2\Vert }=0
\end{equation}
which becomes
\begin{equation}
\gamma _{2}^{2}=\left( 1+\frac{p_{1}}{1+\gamma _{1}}p_{2\Vert }\right) ^{2}
\end{equation}

After squaring and developing both sides,
\begin{equation}
p_{2\Vert }^{2}+p_{2\bot }^{2}=\left( \frac{p_{1}}{1+\gamma _{1}}\right)
^{2}p_{2\Vert }^{2}+\frac{2p_{1}}{1+\gamma _{1}}p_{2\Vert }
\end{equation}
one obtains
\begin{equation}
p_{2\Vert }^{2}\left[ 1-\left( \frac{p_{1}}{1+\gamma _{1}}\right)
^{2}\right] -\frac{2p_{1}}{1+\gamma _{1}}p_{2\Vert }+p_{2\bot }^{2}=0
\end{equation}
and using the relation 
\begin{equation}
1-\left( \frac{p_{1}}{1+\gamma _{1}}\right) ^{2}=\frac{2}{\gamma _{1}+1}
\end{equation}
it turns out that
\begin{equation}
\frac{2p_{2\Vert }^{2}}{\gamma _{1}+1}-\frac{2p_{1}}{\gamma _{1}+1}p_{2\Vert
}+p_{2\bot }^{2}=0
\end{equation}
or 
\begin{equation}
\frac{p_{2\Vert }^{2}}{\gamma _{1}^{2}-1}-\frac{p_{1}}{\gamma _{1}^{2}-1}%
p_{2\Vert }+\frac{p_{2\bot }^{2}}{2\left( \gamma _{1}-1\right) }=0
\end{equation}
which is equivalent to ellipse equation
\begin{equation}
\left( \frac{p_{2\Vert }}{\sqrt{\gamma _{1}^{2}-1}}-\frac{1}{2}\right) ^{2}+%
\frac{p_{2\bot }^{2}}{2\left( \gamma _{1}-1\right) }=\frac{1}{4}
\end{equation}
since
\begin{equation}
\frac{p_{1}}{\gamma _{1}^{2}-1}=\frac{1}{p_{1}}=\frac{1}{\sqrt{\gamma
_{1}^{2}-1}}
\end{equation}

\subsection{Secondary electron source term}

The secondary electron source term, or the scattering rate is given by the M%
\o ller differential relativistic electron-electron cross-section \cite
{hei57}, which has the following form
\begin{equation}
d\sigma _{M}=2\pi r_{e}^{2}\sum \left( \gamma _{1},\gamma _{2}\right)
d\gamma _{2}
\end{equation}
where $r_{e}$ is the classical radius of the electron, and
\begin{eqnarray}
\sum \left( \gamma _{1},\gamma _{2}\right)  &=&\frac{\gamma _{1}^{2}}{\left(
\gamma _{1}^{2}-1\right) \left( \gamma _{2}-1\right) ^{2}\left( \gamma
_{1}-\gamma _{2}\right) ^{2}}\times   \nonumber \\
&&\left\{ \left( \gamma _{1}-1\right) ^{2}-\frac{\left( \gamma _{2}-1\right)
\left( \gamma _{1}-\gamma _{2}\right) }{\gamma _{1}^{2}}\times \left[
2\gamma _{1}^{2}-2\gamma _{1}-1-\left( \gamma _{2}-1\right) \left( \gamma
_{1}-\gamma _{2}\right) \right] \right\} 
\end{eqnarray}
which becomes for very energetic runaways, 
\begin{equation}
\lim_{\gamma _{1}\rightarrow \infty }\sum \left( \gamma _{1},\gamma
_{2}\right) =\frac{1}{\left( \gamma _{2}-1\right) ^{2}}=\frac{1}{E_{c2}^{2}}
\end{equation}
if one introduces the kinetic energy $E_{c2}=\gamma _{2}-1$ of the scattered
electron. It may be expressed also in a simplified form
\begin{equation}
\sum \left( \gamma _{1},\gamma _{2}\right) =\frac{1}{\left( \gamma
_{1}^{2}-1\right) }\left( \frac{\gamma _{1}^{2}\left( \gamma _{1}-1\right)
^{2}}{E_{c2}^{2}\left( \gamma _{1}-1-E_{c2}\right) ^{2}}-\frac{2\gamma
_{1}^{2}+2\gamma _{1}+1}{E_{c2}\left( \gamma _{1}-1-E_{c2}\right) }+1\right) 
\end{equation}
where $E_{c1}=\gamma _{1}-1$ is simply the kinetic energy of the primary
electron after the collision. 

It is now important to separe two regions in the phase space for the primary
electrons, those who belong to the region $\Pi _{i}$ inside the integration
domain of the kinetic equation, and those who have gained a very large
kinetic energy, who belong to the region $\Pi _{o}$ beyond the upper limit
of the integration domain. 

Recalling that $\gamma _{2}=\sqrt{1+p_{2}^{2}},$ $d\gamma _{2}=\left(
p_{2}/\gamma _{2}\right) dp_{2}=v_{2}dp_{2}$ and
\begin{equation}
d\sigma _{M}=2\pi r_{e}^{2}\sum \left( \gamma _{1},\gamma _{2}\right) \left( 
\frac{p_{2}}{\gamma _{2}}\right) dp_{2}=2\pi r_{e}^{2}\sum \left( \gamma
_{1},\gamma _{2}\right) v_{2}dp_{2}
\end{equation}

\subsubsection{ Source term from very energetic runaways}

For very energetic runaway electrons belonging to the region $\Pi _{o},$ the
energy transfer is very small, i.e. 
\begin{equation}
E_{c2}\ll E_{c1}
\end{equation}
and simplifications may be done
\begin{equation}
\sum \left( \gamma _{1},\gamma _{2}\right) \approx \frac{1}{\left( \gamma
_{1}^{2}-1\right) }\left( \frac{\gamma _{1}^{2}}{E_{c2}^{2}}-\frac{2\gamma
_{1}^{2}+2\gamma _{1}+1}{E_{c2}\left( \gamma _{1}-1\right) }+1\right) 
\approx \frac{\gamma _{1}^{2}}{\left( \gamma _{1}^{2}-1\right) }\frac{1}{%
E_{c2}^{2}}
\end{equation}
which becomes in the limit $\gamma _{1}\gg 1,$ 
\begin{equation}
\sum \left( \gamma _{1},\gamma _{2}\right) \approx \frac{1}{E_{c2}^{2}}
\end{equation}

Since $d\gamma _{2}=dE_{c2},$ 
\begin{eqnarray}
d\sigma _{M} &\approx &2\pi r_{e}^{2}\frac{dE_{c2}}{E_{c2}^{2}}  \nonumber \\
&=&-2\pi r_{e}^{2}d\left( \frac{1}{E_{c2}}\right)   \nonumber \\
&=&-2\pi r_{e}^{2}\frac{d\left( 1/E_{c2}\right) }{dp_{2}}dp_{2}  \nonumber \\
&=&2\pi r_{e}^{2}\frac{d}{dp_{2}}\left( \frac{1}{1-\gamma _{2}}\right) dp_{2}
\nonumber \\
&=&2\pi r_{e}^{2}\frac{d}{dp_{2}}\left( \frac{1}{1-\sqrt{1+p_{2}^{2}}}
\right) dp_{2}
\end{eqnarray}

As each $p_{2}$ is uniquely associated with $\xi ^{*},$ and since all
primary electrons are consider to have approximately zero pitch-angle, we
can write
\begin{equation}
d\sigma _{M}\simeq 2\pi r_{e}^{2}\frac{d}{dp_{2}}\left( \frac{1}{1-\sqrt{%
1+p_{2}^{2}}}\right) \delta \left( \xi _{2}-\xi ^{*}\right) d\xi _{2}dp_{2}
\end{equation}
and from this expression, it is possible to derive the the secondary runaway
electron source term $\mathcal{S}_{R}\left( t,\theta ,\psi ,p_{2},\xi
_{2}\right) $ at the poloidal position $\theta ,$ on the magnetic flux
surface $\psi $ and time $t$ for a beam of monoenergetic primary electrons
of velocity $v_{1}$ from the standard definition
\begin{equation}
n_{e}\left( t,\psi \right) n_{R}\left( t,\psi \right) v_{1}d\sigma _{M}=2\pi %
p_{2}^{2}dp_{2}d\xi _{2}\mathcal{S}_{R}\left( t,\theta ,\psi ,p_{2},\xi
_{2}\right) 
\end{equation}
so that
\begin{equation}
\mathcal{S}_{R}\left( t,\theta ,\psi ,p_{2},\xi _{2}\right) =\frac{%
r_{e}^{2}n_{e}\left( t,\psi \right) n_{R}\left( t,\psi \right) v_{1}}{%
p_{2}^{2}}\frac{d}{dp_{2}}\left( \frac{1}{1-\sqrt{1+p_{2}^{2}}}\right)
\delta \left( \xi _{2}-\xi ^{*}\right) 
\end{equation}
where $n_{R}\left( t,\psi \right) $ is the density of primary runaway
electrons and $n_{e}\left( t,\psi \right) $ is the density of target
electrons which are scattered at momentum $p_{2}$ with a pitch-angle cosine $%
\xi ^{*}$. \strut Here $n_{R}$ is only a function of the radial position,
since fast electrons are only passing particles with no magnetic moments
(zero pitch-angle).

Since in the high velocity limit, $v_{1}\approx c,$ one obtains finally 

\begin{equation}
\mathcal{S}_{R}\left( t,\theta ,\psi ,p_{2},\xi _{2}\right) =\frac{%
cr_{e}^{2}n_{e}n_{R}}{p_{2}^{2}}\frac{d}{dp_{2}}\left( \frac{1}{1-\sqrt{%
1+p_{2}^{2}}}\right) \delta \left( \xi _{2}-\xi ^{*}\right) 
\end{equation}

Therefore, dropping the index of the secondary electron which has here no
more interest, the secondary electron source term arising from very
energetic runaway electrons lying in the region $\Pi _{o}$ beyond the upper
integration limit $p_{\max },$ 
\begin{equation}
\mathcal{S}_{R}^{\Pi _{o}}\left( t,\theta ,\psi ,p,\xi \right) =\frac{%
cr_{e}^{2}n_{e}n_{R}^{\Pi _{o}}}{p^{2}}\frac{d}{dp}\left( \frac{1}{1-\sqrt{%
1+p^{2}}}\right) \delta \left( \xi -\xi ^{*}\right) 
\end{equation}
where $\xi ^{*}$ is given by relation \ref{pitchanglerunaway}. This is
exactly the expression given in Ref. \cite{ros97}, where the Coulomb
logarithm comes in through the normalization to the reference collision
frequency $\nu _{e}^{\dag }$ introduced in Sec. \ref{normalization}.
Recalling that
\begin{equation}
\mathcal{S}_{R}^{\Pi _{o}}\left( t,\theta ,\psi ,p,\xi \right) =\frac{%
cr_{e}^{2}n_{e}^{\dag 2}}{\beta _{th}^{\dag 3}}\frac{\overline{n}_{e}%
\overline{n}_{R}^{\Pi _{o}}}{\overline{p}^{2}}\frac{d}{d\overline{p}}\left( 
\frac{1}{1-\sqrt{1+\beta _{th}^{\dag 2}\overline{p}^{2}}}\right) \delta
\left( \xi -\xi ^{*}\right) 
\end{equation}
according to the normalization procedure given in Sec. \ref{normalization},
and using the expression of $\nu _{e}^{\dag },$ one obtains readily
\begin{equation}
\mathcal{S}_{R}^{\Pi _{o}}\left( t,\theta ,\psi ,p,\xi \right) =\frac{
n_{e}^{\dag }\nu _{e}^{\dag }}{\ln \Lambda _{ee}^{\dag }}\frac{\overline{n}%
_{e}\overline{n}_{R}^{\Pi _{o}}}{\overline{p}^{2}}\frac{d}{d\overline{p}}
\left( \frac{1}{1-\sqrt{1+\beta _{th}^{\dag 2}\overline{p}^{2}}}\right)
\delta \left( \xi -\xi ^{*}\right) 
\end{equation}
so that $\overline{\mathcal{S}}_{R}^{\Pi _{o}}\left( t,\theta ,\psi ,p,\xi
\right) =\mathcal{S}_{R}^{\Pi _{o}}\left( t,\theta ,\psi ,p,\xi \right)
/n_{e}^{\dag }\nu _{e}^{\dag },$ the final expression in normalized units
used by the code is
\begin{equation}
\overline{\mathcal{S}}_{R}^{\Pi _{o}}\left( t,\theta ,\psi ,p,\xi \right) =%
\frac{1}{\ln \Lambda _{ee}^{\dag }}\frac{\overline{n}_{e}\overline{n}_{R}^{%
\Pi _{o}}}{\overline{p}^{2}}\frac{d}{d\overline{p}}\left( \frac{1}{1-\sqrt{%
1+\beta _{th}^{\dag 2}\overline{p}^{2}}}\right) \delta \left( \xi -\xi
^{*}\right) 
\end{equation}
which must be bounce-averaged in order to incorporate particle trajectory
effect as done in Sec. \ref{moments}, and remove the poloidal dependence $%
\theta .$

If assumptions of small energy exchange and high-velocity limit for primary
electrons are removed, still assuming a beam of primary electron of velocity 
$v_{1}$ along the magnetic field line direction with no pitch-angle, the
general expression for $\mathcal{S}_{R}\left( t,\theta ,\psi ,p_{2},\xi
_{2}\right) $ becomes
\begin{equation}
\mathcal{S}_{R}\left( t,\theta ,\psi ,p_{2},\xi _{2}\right)
=cr_{e}^{2}n_{e}n_{R}\frac{p_{1}}{\gamma _{1}}\frac{\sum \left( \gamma
_{1},\gamma _{2}\right) }{\gamma _{2}p_{2}}\delta \left( \xi _{2}-\xi
^{*}\right) 
\end{equation}
using the differential relation $d\gamma _{2}=\left( p_{2}/\gamma
_{2}\right) dp_{2}$ and $v_{1}=cp_{1}/\gamma _{1}$ in relativistic units.
After dropping the index for secondary electrons,
\begin{equation}
\mathcal{S}_{R}^{\Pi _{o}}\left( t,\theta ,\psi ,p,\xi \right)
=cr_{e}^{2}n_{e}n_{R}^{\Pi _{o}}\frac{p_{1}}{\gamma _{1}}\frac{\sum \left(
\gamma _{1},\gamma \right) }{\gamma p}\delta \left( \xi -\xi ^{*}\right) 
\end{equation}
and in normalized units for kinetic calculations (******* verifier
normalization********)
\begin{equation}
\overline{\mathcal{S}}_{R}^{\Pi _{o}}\left( t,\theta ,\psi ,p,\xi \right) =%
\overline{n}_{e}\overline{n}_{R}^{\Pi _{o}}\frac{\beta _{th}^{\dag 3}}{\ln
\Lambda _{ee}^{\dag }}\frac{\overline{p}_{1}}{\overline{p}}\frac{\sum \left(
\gamma _{1},\gamma \right) }{\gamma \gamma _{1}}\delta \left( \xi -\xi
^{*}\right)   \label{beamsourcegeneral}
\end{equation}

For a beam of fast primary runaway electrons with no pitch-angle with
various velocities, the distribution function may be expressed in the simple
form $f_{0R}(t,\psi ,\overline{p}_{1},\xi _{1})=f_{R}\left( t,\psi ,%
\overline{p}_{1}\right) \delta \left( \xi _{1}-1\right) $. It is important
to note that $f_{0R}$ has no poloidal dependence, since runaway electrons
are supposed to have a negligible magentic moment. The extrapolation from
formula \ref{beamsourcegeneral} is then straightforward, using the
prescription
\begin{equation}
\overline{n}_{R}^{\Pi _{o}}\left( t,\psi \right) \frac{\overline{p}_{1}}{%
\gamma _{1}}\mathcal{A}\left( \overline{p}_{1}\right) \rightarrow 2\pi \iint 
\mathcal{A}\left( \overline{p}_{1}\right) f_{R}\left( t,\psi ,\overline{p}%
_{1}\right) \delta \left( \xi _{1}-1\right) \frac{\overline{p}_{1}}{\gamma
_{1}}\overline{p}_{1}^{2}d\overline{p}_{1}d\xi _{1}=2\pi \int_{p_{\max }}^{%
\infty }\mathcal{A}\left( \overline{p}_{1}\right) f_{R}\left( \overline{p}%
_{1}\right) \frac{\overline{p}_{1}^{3}}{\gamma _{1}}d\overline{p}_{1}
\end{equation}
where $\mathcal{A}\left( \overline{p}_{1}\right) $ is a function of $%
\overline{p}_{1}$ arising from the energy dependence of the e-e collision
cross-section. Here, the lower integral boundary is taken to be the upper
limit of the Fokker-Planck integration domain, since in this Section, we
consider only the knock-on process resulting from very energetic primary
electrons whose momentum is well beyond $p_{\max }$.

\subsubsection{Source term from moderate energetic runaways}

In that case, it is no more possible to neglect the momentum distribution
function $f_{0}$ of primary electrons which have kinetic energies $E_{c1}$
larger than $E_{*}$ the critical runaway energy threshold. Inside the region 
$\Pi _{i,}$the full expression of the M\o ller cross-section $d\sigma _{M}$
must be used. The generalized source function is then given by
\[
\mathcal{S}_{R}\left( p_{2},\xi _{2}\right) =cr_{e}^{2}n_{e}\frac{p_{1}}{%
\gamma _{1}}\frac{1}{\gamma _{2}p_{2}}\sum \left( \gamma _{1},\gamma
_{2}\right) f_{0}\left( \psi ,\theta ,p_{1},\xi _{1}\right) \delta \left(
\xi _{2}-\xi ^{*}\right) 
\]

\subsection{Critical electric field}
